
\subsection{Integral and Gaussian descriptions}

\begin{statement}[C.8.6; Proved: the Mahler-measure identities]
Let $\Lambda\subset\C$ be a lattice of area $A$, let $v$ be a zero
of $\wp_\Lambda$, and let $d\mu=A^{-1}dA_z$ be Haar probability
measure on $\C/\Lambda$. Put $U(z)=\log N(z)$ and
$L=U(v)$. Then
\begin{align}
 6L&=\int_{\C/\Lambda}\log\left|\frac{\wp'(z)^2}{\wp(z)^3}\right|d\mu(z),
 \label{eq:ca-mahler}\\
 N(v)^2&=|\Delta(\Lambda)|^{1/6}
      \exp\left(-\int_{\C/\Lambda}\log|\wp(z)|d\mu(z)\right).
 \label{eq:ca-wpmean}
\end{align}
If $C(u)=1-(g_2/4)u^2-(g_3/4)u^3$, the pushforward of $d\mu$
under $u=1/\wp(z)$ is
\begin{equation}\label{eq:ca-nu}
 d\nu(u)=\frac1{2A}\frac{dA_u}{|u|\,|C(u)|}.
\end{equation}
In particular its unnormalized integral is exactly $2A$, and
\[
 N(v)^6=4\exp\left(\int_{\C}\log|C(u)|d\nu(u)\right).
\]
Writing $e_1,e_2,e_3$ for the three branch values and
$p(x)=\int\log|x-\wp(z)|d\mu(z)$, the maximum theorem is equivalently
\[
 \int\log|C|d\nu\leq\log2,\qquad
 p(0)\geq\frac{p(e_1)+p(e_2)+p(e_3)}3-\frac{\log2}3.
\]
Equality in either inequality occurs exactly when $g_3=0$.
\end{statement}
\begin{proof}
The mean of $U$ is zero by the theta product calculation in C.1,
and near zero
$U(z)=\log|z|+\frac1{12}\log|\Delta|+o(1)$.
The divisor of $\wp$ is $(v)+(-v)-2(O)$, with multiplicity
two when $v=-v$. Subtract its logarithmic divisor from the
corresponding combination of $U$'s. The result is harmonic
on the compact torus and hence constant. Its expansion at
zero gives
\[
 \log|\wp(z)|=U(z-v)+U(z+v)-2U(z)
                  +\frac16\log|\Delta|-2L.
\]
Integration proves \eqref{eq:ca-wpmean}.
Similarly the divisor of $\wp'$ is the sum of the three
nonzero half-periods minus $3(O)$. The multiplication-by-two
identity of C.1 gives the product of $N$ at those half-periods
as $2$. Since $\wp'(z)\sim-2z^{-3}$, comparison at zero yields
\[
 \log|\wp'(z)|=\sum_{w\in E[2]\setminus\{O\}}U(z-w)-3U(z)
                         +\frac14\log|\Delta|.
\]
Its mean is $\frac14\log|\Delta|$. Twice this identity minus
three times the preceding mean proves \eqref{eq:ca-mahler}.
All logarithmic singularities are locally integrable, so these
integrations also cover a zero at a half-period.

The map $z\mapsto x=\wp(z)$ has degree two. Its two inverse
sheets each have Jacobian $|\wp'(z)|^{-2}$; branch values form
a measure-zero set. Thus its pushforward density is
$2/(A|4x^3-g_2x-g_3|)$. Under $x=1/u$ the Jacobian is
$|u|^{-4}$ and the denominator is $4|u|^{-3}|C(u)|$.
This proves \eqref{eq:ca-nu} and its normalization. The identity
$\wp'^2/\wp^3=4C(u)$ gives the next formula.

Replacing $v$ above by a half-period gives
$p(e_j)=\frac16\log|\Delta|-2U(w_j)$; summing and using
$\sum_jU(w_j)=\log2$ proves the centroid formulation.
The equality characterization is that of C.8.5,
since $j=1728$ is equivalent to $g_3=0$ for a nonsingular lattice.
\end{proof}

\begin{statement}[C.8.7; Proved: the two complex beta integrals]
For the square and equianharmonic lattices respectively, the
integrals in C.8.6 are
\[
 \int\log|C|\,d\nu=\log2,\qquad
 \int\log|C|\,d\nu=\frac32\log3-2\log2.
\]
They recover $N^2=2$ and $N^4=3$ with the same normalization.
\end{statement}
\begin{proof}
By homogeneity it suffices to take $C(u)=1-u^m$, with $m=2$
or $m=3$. For $-1<\Re s<1-1/m$ put
\[
 I_m(s)=\int_{\C} |u|^{-1}|1-u^m|^{s-1}dA_u.
\]
The substitution $w=u^m$, with all $m$ sheets counted, gives
\[
 I_m(s)=\frac1m B_{\C}\left(\frac1{2m},\frac{s+1}2\right),
 \quad
 B_{\C}(a,b)=\pi\frac{\Gamma(a)\Gamma(b)\Gamma(1-a-b)}
                   {\Gamma(1-a)\Gamma(1-b)\Gamma(a+b)}.
\]
We include a proof of the integral evaluation. For
$0<a,b$ and $a+b<1$, insert
\[
 |w|^{2a-2}=\frac1{\Gamma(1-a)}
       \int_0^\infty t^{-a}e^{-t|w|^2}\,dt
\]
and the corresponding formula for $|1-w|^{2b-2}$.
Tonelli's theorem applies to the positive integrand. The
Gaussian integral over $w$ is
$\pi(t+s)^{-1}\exp(-ts/(t+s))$.
Putting $t=rv$, $s=r(1-v)$ reduces the remaining integrals
to $\Gamma(1-a-b)$ and the beta integral $B(a,b)$, proving
the displayed formula. The same integrable bounds justify
differentiation in a neighbourhood of $s=0$.

Consequently, with $a=1/(2m)$ and $\psi=\Gamma'/\Gamma$,
\[
 \frac{I_m'(0)}{I_m(0)}
 =\frac12\left(2\psi(1/2)-\psi(1/2-a)-\psi(1/2+a)\right).
\]
The gamma duplication and triplication formulas, differentiated
logarithmically, evaluate this as $\log2$ for $m=2$ and
$\frac32\log3-2\log2$ for $m=3$. These can also be checked
directly from
$\psi(1/2)=-\gamma-2\log2$,
$\psi(1/4)+\psi(3/4)=-2\gamma-6\log2$ and
$\psi(1/3)+\psi(2/3)=-2\gamma-3\log3$.
Equation \eqref{eq:ca-mahler} completes the normalization test.
\end{proof}

\begin{statement}[C.8.8; Proved: the signed Gaussian identity]
For $\Lambda=\Z+\tau\Z$, $T=\Im\tau$, and every $z\in\C$,
\[
 \sqrt{2T}\,|\eta(\tau)|^2N(z,\tau)^2
 =\sum_{m\in\frac12\Lambda}\epsilon(m)
                         e^{-2\pi|z-m|^2/T},
\]
where $\epsilon=-1$ on $\Lambda$ and $+1$ on the three other
cosets of $\Lambda$ in $\frac12\Lambda$.
Both sides and all their derivatives converge locally uniformly.
\end{statement}
\begin{proof}
Write $z=b+a\tau$ and expand in the Fourier characters
$e^{2\pi i(kb+la)}$ of $(a,b)\in(\R/\Z)^2$.
For $G(z)=\sum_{\lambda\in\Lambda}e^{-2\pi|z-\lambda|^2/T}$,
Poisson summation gives the coefficient
\[
 \frac12\exp\left(-\frac{\pi|k\tau-l|^2}{2T}\right).
\]
The right side is the sum of the four half-period translates
of $G$, minus $2G$. Its coefficient is therefore
\[
 (-1)^{k+l+kl}
       \exp\left(-\frac{\pi|k\tau-l|^2}{2T}\right).
\]
To check the sign, it is $+1$ when both $k,l$ are even and
$-1$ otherwise.

On the left use
\[
 \vartheta_1(z,\tau)
 =-i\sum_{n\in\Z}(-1)^n
        e^{\pi i(n+1/2)^2\tau+2\pi i(n+1/2)z}.
\]
Multiply by its conjugate and by
$\sqrt{2T}e^{-2\pi(\Im z)^2/T}$.
In the double sum put $k=n-m$ and apply one-dimensional
Poisson summation to the remaining Gaussian in $m$.
Completing the square gives precisely the coefficient just
displayed. This proves equality of the Fourier series.
Positive definiteness of $|k\tau-l|^2$ and Gaussian decay
justify the changes of order, Poisson summation, and every
claimed differentiation.
\end{proof}

